\documentclass[10pt,a4paper]{amsart}
\usepackage[margin=19mm]{geometry}
\usepackage{amsmath,amssymb,mathtools}
\usepackage[scr=rsfs,cal=euler]{mathalfa}
\usepackage{xfrac}
\usepackage[unicode,hidelinks,bookmarks=false]{hyperref}
\newcommand{\ie}{i.e.\ }
\newcommand{\cf}{cf.\ }
\newcommand{\resp}{respectively\ }
\newcommand{\Subsection}{\subsection}
\providecommand{\nobreakdash}{\nobreak\hbox{-}}
\setlength{\emergencystretch}{2em}
\multlinegap0pt
\usepackage{mdframed}
\usepackage{mdframed}
\usepackage{mdframed}
\usepackage{mdframed}
\usepackage{mathtools}
\usepackage{amssymb}
\usepackage{bm}
\usepackage{dsfont}
\usepackage{tikz}
\usepackage{xcolor}
\usepackage{float}
\usepackage[shortlabels]{enumitem}
\usepackage{mdframed}
\newtheorem{lemma}{Lemma}
\title{A Robust Optimization Approach to Sparse Principal Component Analysis}
\author{David V\"avinggren, Francis Bach, Andr\'e M. H. Teixeira, Dave Zachariah, Ant\^onio H. Ribeiro}
\date{Source: arXiv:2606.03553v1 (CC BY 4.0)}
\begin{document}
\maketitle

\noindent\emph{Source and licence.} A Robust Optimization Approach to Sparse Principal Component Analysis. Version arXiv:2606.03553v1 (CC BY 4.0), distributed by arXiv under the Creative Commons Attribution 4.0 International licence, http://creativecommons.org/licenses/by/4.0/, as recorded in the arXiv licence metadata for this exact version and shown on the abstract page https://arxiv.org/abs/2606.03553v1. Full licence notice: \texttt{licenses/arxiv-2606.03553-CC-BY-4.0.txt}.\par\smallskip

\noindent\textit{Editorial scope.} Counted unit: the article's lemma lemma (source0.tex lines 1259--1261). The article's complete original proof of it is delivered with it (source0.tex lines 1262--1282, one \texttt{proof} environment of 21 lines). No mathematics is written, repaired or re-derived: the statement and the proof are the article's own lines, in the article's own notation, and no retained line is substituted or re-flowed.  No further source-local context is needed: every cross-reference of every retained line resolves inside the two delivered spans.  Nothing was omitted from any retained span, so the package carries no omission record.  No retained line contains a citation, so the wrapper carries no bibliography and no bibliography entry.  No retained line contains a figure environment, an image-inclusion command or drawing code, so no image is delivered and the package declares no asset dependency.  Editorial formatting only, in addition: an AMS \texttt{amsart} preamble that restates the article's own declarations and macros used by the retained lines, namely the environments \texttt{lemma} on separate counters, together with the standard packages those lines need.  The retained environments print in this document as Lemma 1 in order of appearance, whereas the article prints the same items with the numbering of its own full document.  The complete native source of the article as distributed in the arXiv e-print for this exact version is bundled unchanged as \texttt{sources/201984/source0.tex}.
\begin{lemma}
    Consider $f: \mathbb{R} \to \mathbb{R}$ given by $f(x) = x^2 - bx - c$ and $g: \mathbb{R} \to \mathbb{R}$ given by $g(x) = x^2 - b^2 - 2c$ where $b,c \in \mathbb{R}$. If $b^2 + 4c \geq 0$ and $f(x) \leq 0$, then $g(x) \leq f(x)$.
\end{lemma}

\begin{proof}
    We aim to show that $f(x) - g(x) \geq 0$. First, we expand the difference between the two functions:
    \[
        f(x) - g(x) = (x^2 - bx - c) - (x^2 - b^2 - 2c) = b^2 - bx + c.
    \]
    Observe that we can construct this exact expression by expanding the perfect square $(x-b)^2$ and subtracting our function $f(x)$:
    \[
        (x-b)^2 - f(x) = (x^2 - 2bx + b^2) - (x^2 - bx - c) = b^2 - bx + c.
    \]
    Therefore, we establish the direct algebraic identity:
    \[
        f(x) - g(x) = (x-b)^2 - f(x).
    \]
    Since the square of any real number is non-negative, we know $(x-b)^2 \geq 0$. We are explicitly given the condition that $f(x) \leq 0$, which implies $-f(x) \geq 0$. 
    Because $f(x) - g(x)$ is expressed as the sum of two non-negative terms, it must also be non-negative:
    \[
        f(x) - g(x) \geq 0 \implies g(x) \leq f(x).
    \]
    \textit{Note:} The condition $b^2 + 4c \geq 0$ guarantees that the premise $f(x) \leq 0$ is attainable for some $x \in \mathbb{R}$.

\end{proof}

\end{document}
